%==============================================================================
%  WEEKLY AI BRIEF -- issue #06, week ending 29 September 2026
%  Engine: tectonic
%==============================================================================
\documentclass[11pt,a4paper]{article}

%------------------------------ packages -------------------------------------
\usepackage[T1]{fontenc}
\usepackage[utf8]{inputenc}
\usepackage{lmodern}
\usepackage[margin=2.2cm,top=2.0cm,bottom=2.3cm]{geometry}
\usepackage{microtype}
\usepackage{graphicx}
\usepackage{xcolor}
\usepackage{titlesec}
\usepackage{enumitem}
\usepackage{booktabs}
\usepackage{tabularx}
\usepackage{fancyhdr}
\usepackage{lastpage}
\usepackage{tcolorbox}
\usepackage[hidelinks]{hyperref}

\tcbuselibrary{skins,breakable}

\graphicspath{{../../../template/}{./}{template/}}

%------------------------------ palette --------------------------------------
\definecolor{brandink}{HTML}{103868}    % KAUST navy -- headings, rules
\definecolor{brandaccent}{HTML}{0B7480} % KAUST teal (darkened to 5.6:1 on white)
\definecolor{brandsoft}{HTML}{EAF3F5}   % box fills
\definecolor{brandgrey}{HTML}{5A6672}   % meta text

\hypersetup{colorlinks=true,linkcolor=brandaccent,urlcolor=brandaccent,
            citecolor=brandaccent,pdfborder={0 0 0}}

%------------------------------ metadata -------------------------------------
\newcommand{\briefnumber}{06}
\newcommand{\briefweek}{23--29 September 2026}
\newcommand{\briefdate}{29 September 2026}
\newcommand{\briefauthor}{Ali Habibullah}
\newcommand{\briefaffil}{KAUST}
\newcommand{\briefscope}{Papers, model releases, benchmarks and policy from the past seven days}

%------------------------------ section styling ------------------------------
\titleformat{\section}
  {\normalfont\sffamily\large\bfseries\color{brandink}}{\thesection.}{0.6em}{}
  [\vspace{-0.55em}{\color{brandaccent}\rule{\textwidth}{0.9pt}}]
\titlespacing*{\section}{0pt}{1.5ex plus .3ex}{0.9ex}

\titleformat{\subsection}[runin]
  {\normalfont\sffamily\bfseries\color{brandink}}{}{0pt}{}[\;---]
\titlespacing*{\subsection}{0pt}{1.25ex plus .2ex}{0.5em}

\setlist[itemize]{leftmargin=1.15em,itemsep=0.25em,topsep=0.35em,parsep=0pt}
\setlength{\parindent}{0pt}
\setlength{\parskip}{0.45em}

%------------------------------ header / footer ------------------------------
\pagestyle{fancy}
\fancyhf{}
\renewcommand{\headrulewidth}{0pt}
\renewcommand{\footrulewidth}{0.4pt}
\fancyfoot[L]{\footnotesize\color{brandgrey}Weekly AI Brief \#\briefnumber\ \textbullet\ \briefweek}
\fancyfoot[R]{\footnotesize\color{brandgrey}Page \thepage\ of \pageref{LastPage}}

%------------------------------ content macros -------------------------------
\newcommand{\tldr}[1]{%
  \begin{tcolorbox}[breakable,enhanced,colback=brandsoft,colframe=brandaccent,
      boxrule=0.8pt,left=10pt,right=10pt,top=8pt,bottom=8pt,arc=2pt,
      title={\sffamily\bfseries The week in five lines},fonttitle=\color{white}]
    \begin{itemize}[leftmargin=1.1em,itemsep=0.3em,topsep=0pt]#1\end{itemize}
  \end{tcolorbox}}

\newcommand{\entry}[4]{%
  \par\vspace{0.7em}%
  {\sffamily\bfseries\color{brandink}#1}\par
  {\footnotesize\color{brandgrey}#2 \textbullet\ \href{#3}{link} \textbullet\ \srcref{#4}}\par
  \vspace{0.15em}}

\newcommand{\why}[1]{\par{\small\textbf{Why it matters.} #1}\par}

\newcommand{\srcref}[1]{\textcolor{brandaccent}{[S#1]}}

\newcommand{\stat}[2]{{\sffamily\bfseries\color{brandink}#1}~{\footnotesize\color{brandgrey}#2}}

%==============================================================================
\begin{document}

%------------------------------ title block ----------------------------------
\noindent
\begin{minipage}[c]{0.48\textwidth}
  \raggedright\includegraphics[height=1.4cm]{KAUST_Academy_Logo.png}
\end{minipage}\hfill
\begin{minipage}[c]{0.48\textwidth}
  \raggedleft\includegraphics[height=1.5cm]{KAUST_Logo.png}
\end{minipage}
\vspace{0.7em}

{\sffamily
 {\color{brandgrey}\footnotesize ISSUE \#\briefnumber\ \textbullet\ \briefdate}\par
 \vspace{0.15em}
 {\LARGE\bfseries\color{brandink}Weekly AI Brief}\par
 \vspace{0.25em}
 {\large\color{brandaccent}\briefweek}\par
 \vspace{0.35em}
 {\footnotesize\color{brandgrey}\briefscope\\
  Prepared by \briefauthor, \briefaffil}\par}
\vspace{0.3em}
{\color{brandink}\rule{\textwidth}{1.4pt}}
\vspace{0.6em}

%------------------------------ 1. executive summary -------------------------
\tldr{
  \item \textbf{Claude Sonnet 5.5 ships at \$2/\$10 per million tokens.} Artificial Analysis independently places it 3rd of 216 models, two points behind Opus 5.5. \srcref{1}\,\srcref{2}
  \item \textbf{OpenAI halved its mid-tier prices.} GPT-6 Sol and Luna launched on 22 September at 50\% below GPT-5.6 pricing, about 90 minutes after Anthropic's Opus 5.5. \srcref{3}\,\srcref{5}\,\srcref{19}
  \item \textbf{Xiaomi posted MIT-licensed 1.02T-parameter MiMo-V2.6 checkpoints.} The 27 September update targets agents that repeat the same tool call. \srcref{7}
  \item \textbf{Claude agents flagged a new CRISPR-like enzyme family in phage DNA.} Anthropic's own preprint says its function is still untested. \srcref{9}\,\srcref{10}
  \item \textbf{The UN Security Council heard from Altman, Amodei and Bengio.} The US told the Council that dialogue ``cannot be allowed to drift towards global governance''. \srcref{15}
}

%------------------------------ 2. models & releases -------------------------
\section{Model and product releases}

\entry{Claude Sonnet 5.5: near-Opus scores at half the Opus price}{Anthropic \textbullet\ API and all major clouds \textbullet\ 28 September}{https://www.anthropic.com/claude-sonnet-5-5}{1}
Sonnet 5.5 (\texttt{claude-sonnet-5-5}) costs \$2 input and \$10 output per million tokens,
with cache reads at \$0.20. It is available on AWS, Google Cloud and Microsoft Azure.
Anthropic says it is 30\%+ faster than Sonnet~5 and costs up to 30\% less per task.
Its largest self-reported jump is on Terminal-Bench~4.0, from 10.3\% to 70.6\%, which
is above Opus~5.5's 66.4\% on the same page. Artificial Analysis independently scores it
56 on Intelligence Index v4.3.2 (\#3 of 216). It also lists a 1M-token context window
and flags the model as ``very verbose'': 410M output tokens across the evaluation,
against an 88M median. \srcref{2}
\why{For most agentic coding workloads, the default question is now whether Opus is worth twice the price, not whether Sonnet is good enough.}

\entry{GPT-6 Sol and Luna: OpenAI halves prices below the GPT-6 flagship}{OpenAI \textbullet\ API, ChatGPT Work, Codex \textbullet\ 22 September (one day before this window)}{https://openai.com/index/introducing-gpt-6-sol-and-luna/}{3}
Carried here because they launched on the last day of issue \#05's window, after it was compiled.
Sol drops from \$4/\$20 to \$2/\$10 and Luna from \$0.20/\$1.20 to \$0.10/\$0.50 per million tokens.
OpenAI says both models were trained ``with similar methods as GPT-6 Astra''.
It self-reports DeepSWE~v1.1 at 68.8\% for Sol (max effort), 1.1 points below Claude Fable~5
at ``approximately 80\% lower cost per task''. On its internal factuality evaluation,
Sol ``makes about half as many mistakes'' as its predecessor. OpenAI notes that this
evaluation is built from conversations where users flagged errors, so it is not representative of
typical use. Cached input reads are discounted 90\%. The API IDs are \texttt{gpt-6-sol} and
\texttt{gpt-6-luna}. Artificial Analysis scores Sol (max) at 48 (\#20 of 216) with an 872K context window. \srcref{4}
\why{Bulk extraction and summarisation jobs that were routed to small models on cost grounds should be re-priced against Luna at \$0.10/\$0.50.}

\entry{Claude Opus 5.5 tops the independent index}{Anthropic \textbullet\ API and all major clouds \textbullet\ 22 September (one day before this window)}{https://www.anthropic.com/claude-opus-5-5}{5}
Also a 22 September launch that issue \#05 missed. Opus~5.5 costs \$4/\$20 per million tokens
(\$8/\$40 in fast mode), which Anthropic says is ``40\% less to run than Opus~5'' on
typical workloads. Its self-reported table puts it ahead of GPT-6 Astra on Terminal-Bench~4.0
(66.4\% vs 57.9\%) and behind Astra on AutomationBench (40.0\% vs 41.4\%) and Terminal-Bench-Science
(58.7\% vs 64.6\%). Artificial Analysis independently ranks it \#1 of 216 on Intelligence Index
v4.3.2, at 58. \srcref{6}
\why{The model at the top of the independent index is also, by Anthropic's account, 40\% cheaper to run than Opus~5, so teams that stayed on Fable~5.1 for quality should re-run their evals on Opus~5.5.}

\entry{MiMo-V2.6 open weights, 1.02T parameters, MIT licence}{Xiaomi \textbullet\ Hugging Face and ModelScope \textbullet\ MOPD checkpoints posted 27 September}{https://huggingface.co/XiaomiMiMo/MiMo-V2.6-Pro-MOPD}{7}
Pro is a sparse MoE with 1.02T total and 42B active parameters. Flash has 309B total
and 15B active. Both have a 1M-token context and native text, image, video and audio input.
The RL checkpoints appeared on 21--22 September. The in-window 27 September ``MOPD''
checkpoints add multi-teacher on-policy distillation and target tool-call repetition,
a failure mode where an agent ``keeps issuing the same or highly similar tool calls,
appearing busy while making no progress''. Xiaomi's self-reported table for the RL checkpoint puts Pro at
71.9 on DeepSWE~v1.1, against 74.0 for Claude Opus~5. \srcref{8}
\why{A trillion-parameter checkpoint under MIT removes the licence question for fine-tuning and commercial use, and the repetition failure mode it names is worth checking for in any agent harness.}

%------------------------------ 3. research ----------------------------------
\section{Research worth reading}

\entry{Claude agents flag a CRISPR-like reverse-transcriptase family}{Yoon, Athukoralage, Ameisen et al., Anthropic \textbullet\ preprint, 23 September}{https://www.anthropic.com/news/claude-discovers-novel-enzyme-system}{9}
Claude Code agents surveyed reverse-transcriptase loci across 1.9 billion protein clusters.
They surfaced ``array-associated RTs'' (ART): an RT, a dedicated partner gene and a long
array of roughly 200-nucleotide repeat units, found in bacteriophage DNA. The preprint
reports 95 distinct ART RT clusters, 28 of them with detectable upstream arrays. \srcref{10}
Anthropic's post puts the campaign at about 950 agents, 21 hours and 210 million tokens.
The preprint is explicit that ``ART's biochemical activity and biological function remain
untested''. What the agents did was anomaly detection, not a demonstrated mechanism.
\why{It is a concrete, checkable case of an agent campaign surfacing an uncharacterised enzyme family; whether it matters biologically depends on functional assays the preprint says have not been done.}

\entry{Just-in-Time Memory: curate agent memory at read time, not write time}{Zhou, Li, Liu, Yavuz, Joty \textbullet\ arXiv:2609.27334 \textbullet\ 23 September}{https://arxiv.org/abs/2609.27334}{11}
Most agent memory systems compress each finished trajectory into a fixed reflection or
skill before anyone knows what future task will need it. JitMem keeps raw trajectories.
A trained curator writes a task-specific payload only when a new task arrives, which also
lets it learn from immediate task success. Over the strongest baseline it gains 16.2
(ALFWorld), 16.3 (WebShop) and 3.9 ($\tau^2$-bench) absolute success-rate points. The
arXiv page links no code.
\why{If you are building agent memory, storing raw traces and summarising on retrieval is cheap to try and beat learned write-time summaries here.}

\entry{A measurable definition of ``interesting'' theorems}{Patel, Rammal, Hayat, Munos, Kempe \textbullet\ arXiv:2609.28603 \textbullet\ 23 September}{https://arxiv.org/abs/2609.28603}{12}
The authors define a theorem's interestingness as the ratio of proof length to statement
length. They train a 27B model to predict proof difficulty and use it to select
LLM-generated theorems for a growing machine-verified library. Selection cuts overlap with
existing Mathlib from 91.9\% to 30.6\%. No code link is given on the arXiv page.
\why{Automated maths has had a verifier but no filter; a cheap, checkable proxy for novelty is the missing piece for self-expanding formal libraries.}

\entry{Transformers can hold two next-token distributions at once}{Tikhonov, Korznikov, \ldots, Oseledets, Tutubalina \textbullet\ arXiv:2609.29845 \textbullet\ 24 September}{https://arxiv.org/abs/2609.29845}{13}
When inputs from two text streams are linearly combined, the model outputs a superposition
of the two next-token distributions. The authors find this property is architectural rather
than learned, and that training weakens it. Fine-tuning restores it, and a new decoding
method produces two coherent continuations from a single forward pass. No code is linked.
\why{If this reproduces at scale, it is a new route to batched or speculative generation that needs no extra model.}

\entry{PISA: block-sparse attention in \texorpdfstring{$O(N\log N)$}{O(N log N)}}{Tang, Qin, Pan, Li, Liu \textbullet\ arXiv:2609.31093 \textbullet\ 25 September}{https://arxiv.org/abs/2609.31093}{14}
Block-sparse attention still spends quadratic time choosing which blocks to keep. PISA
instead picks blocks by a pyramid Top-$K$ search over $O(\log N)$ pooled key levels, using
fused Triton kernels for training and inference. Against its baseline it is comparable
on commonsense reasoning and better on retrieval. The abstract gives no speed numbers
and no code link, so the wall-clock gain is unverified.
\why{Block selection is the remaining quadratic step in long-context sparse attention; this removes it, if the kernels hold up.}

%------------------------------ 4. benchmarks & evals ------------------------
\section{Benchmarks, evaluations and reproductions}

{\small
\begin{tabularx}{\textwidth}{@{}X r r r l@{}}
\toprule
\textbf{Benchmark} & \textbf{Sonnet 5.5} & \textbf{Opus 5.5} & \textbf{GPT-6 Sol} & \textbf{Reported by} \\
\midrule
AA Intelligence Index v4.3.2 & 56 & 58 & 48 & Artificial Analysis \\
FrontierCode 1.1 (Main)      & 46.2\% & 54.4\% & 49.3\% & Anthropic \\
GDPval-AA v2.1 (Elo)         & 1844 & 1846 & 1487 & Anthropic \\
AA-Briefcase v1.1 (Elo)      & 1811 & 1822 & 1483 & Anthropic \\
Chartography, no tools       & 61.6\% & 64.4\% & 53.6\% & Anthropic \\
Terminal-Bench 4.0           & 70.6\% & 66.4\% & --- & Anthropic \\
OSWorld 2.1                  & 80.1\% & 81.8\% & --- & Anthropic \\
\bottomrule
\end{tabularx}}
\vspace{0.3em}

{\footnotesize\color{brandgrey}Sources: row 1 \srcref{2}\,\srcref{6}\,\srcref{4}; rows 2--7 from Anthropic's Sonnet~5.5 table \srcref{1}. Only row~1 is an independent measurement.}

Only the first row is independent. On it, Opus~5.5 and Sonnet~5.5 sit 10 and 8 points
above GPT-6 Sol (max). The remaining rows are Anthropic's own and include its choice of
harness for the OpenAI column. Two harness gaps showed up this week. First, Chartography for
Opus~5.5 is 64.4\% ``no tools'' on the Sonnet page and 89.0\% ``with tools'' on the Opus
page \srcref{5}. The benchmark and model are the same; only tool access differs.
Second, Claude Opus~5 on AutomationBench~v1.0.6 is 26.9\% in both OpenAI's and
Anthropic's tables \srcref{3}\,\srcref{5}, but 50.3 in Xiaomi's \srcref{8}. Cross-lab
comparisons on agent benchmarks are not reliable until labs publish their harness settings.

%------------------------------ 5. industry & policy -------------------------
\section{Industry, funding and policy}
\begin{itemize}
  \item \textbf{UN Security Council, 10228th meeting (23 September).} The meeting was chaired by France,
  with Bengio, Altman, Amodei and Hugging Face's Delangue as briefers. Amodei proposed ``a ban on using AI to
  make biological weapons'', verification systems, and common testing standards with ``a notification
  system for AI incidents''. The White House's Michael Kratsios answered that international dialogue
  ``cannot be allowed to drift towards global governance''. \srcref{15}
  \item \textbf{Ban Artificial Superintelligence Act (introduced 23 September).} Sen.\ Sanders and Rep.\ Casar's
  bill would ban AI that ``exceeds human cognitive performance and capabilities across most domains'',
  pause advanced development until a new cabinet-level Department of AI sets rules, and allow up to
  20 years' imprisonment. It is a bill introduced by two members, not law. \srcref{16}
  \item \textbf{Oregon Executive Order 26-26 (23 September).} The order directs the state CIO to draft third-party
  AI-safety review standards and to assess a ``kill switch'' requirement before agencies procure frontier
  models. The CIO's proposal is due within 90 days. It binds state procurement only. \srcref{17}
\end{itemize}

%------------------------------ 6. what to watch -----------------------------
\section{What to watch next week}
\begin{itemize}
  \item \textbf{OpenAI DevDay, San Francisco, 29 September.} The event falls on this issue's
  publication date \srcref{18}. What actually ships (as opposed to the rumoured lists circulating
  beforehand) is next week's lead.
  \item \textbf{Independent scores beyond Artificial Analysis for Sonnet~5.5 and GPT-6 Sol.} The
  vendor tables disagree on harness settings. The first third-party agent-benchmark runs will
  show whether Sonnet~5.5's Terminal-Bench lead over Opus~5.5 holds.
  \item \textbf{Wet-lab follow-up on ART.} The open question is whether any outside lab reports
  expression or activity data for the array-associated RTs, which is what would turn an
  agent-found anomaly into a finding.
\end{itemize}

%------------------------------ 7. sources -----------------------------------
\clearpage
\section{Sources}
{\footnotesize
\begin{tabularx}{\textwidth}{@{}l X@{}}
\toprule
\textbf{\#} & \textbf{Reference} \\
\midrule
S1 & Anthropic. \emph{Introducing Claude Sonnet 5.5}. 28 Sep 2026. \url{https://www.anthropic.com/claude-sonnet-5-5} \\
S2 & Artificial Analysis. \emph{Claude Sonnet 5.5}. Live page, accessed 29 Sep 2026. \url{https://artificialanalysis.ai/models/claude-sonnet-5-5} \\
S3 & OpenAI. \emph{Introducing GPT-6 Sol and Luna}. 22 Sep 2026. \url{https://openai.com/index/introducing-gpt-6-sol-and-luna/} \\
S4 & Artificial Analysis. \emph{GPT-6 Sol (max)}. Live page, accessed 29 Sep 2026. \url{https://artificialanalysis.ai/models/gpt-6-sol} \\
S5 & Anthropic. \emph{Introducing Claude Opus 5.5}. 22 Sep 2026. \url{https://www.anthropic.com/claude-opus-5-5} \\
S6 & Artificial Analysis. \emph{Claude Opus 5.5}. Live page, accessed 29 Sep 2026. \url{https://artificialanalysis.ai/models/claude-opus-5-5} \\
S7 & Xiaomi MiMo. \emph{MiMo-V2.6-Pro-MOPD model card}. 27 Sep 2026. \url{https://huggingface.co/XiaomiMiMo/MiMo-V2.6-Pro-MOPD} \\
S8 & Xiaomi MiMo. \emph{MiMo-V2.6-Pro-RL model card}. 21 Sep 2026. \url{https://huggingface.co/XiaomiMiMo/MiMo-V2.6-Pro-RL} \\
S9 & Anthropic. \emph{Claude discovers a novel enzyme system with CRISPR-like repeats}. 23 Sep 2026. \url{https://www.anthropic.com/news/claude-discovers-novel-enzyme-system} \\
S10 & Yoon, P.\,H. et al. \emph{Autonomous AI agents discover reverse transcriptases with tandem repeat arrays}. alphaXiv, 23 Sep 2026. \url{https://www.alphaxiv.org/abs/2609.ai-agents-discover-reverse-transcriptases} \\
S11 & Zhou, Y. et al. \emph{Just-in-Time Memory}. arXiv:2609.27334, 23 Sep 2026. \url{https://arxiv.org/abs/2609.27334} \\
S12 & Patel, N. et al. \emph{Learning to Discover Interesting Mathematics}. arXiv:2609.28603, 23 Sep 2026. \url{https://arxiv.org/abs/2609.28603} \\
S13 & Tikhonov, P. et al. \emph{Your Transformer Can Hold Two Thoughts at Once}. arXiv:2609.29845, 24 Sep 2026. \url{https://arxiv.org/abs/2609.29845} \\
S14 & Tang, B. et al. \emph{Block Sparse Attention with Log-Linear Complexity}. arXiv:2609.31093, 25 Sep 2026. \url{https://arxiv.org/abs/2609.31093} \\
S15 & United Nations. \emph{Security Council, 10228th meeting: AI and international security}. 23 Sep 2026. \url{https://transcripts.un.org/en/sc/10228} \\
S16 & Office of Rep.\ Greg Casar. \emph{Casar, Sanders introduce legislation to create new federal agency, ban artificial superintelligence}. 23 Sep 2026. \url{https://casar.house.gov/media/press-releases/news-casar-sanders-introduce-legislation-create-new-federal-agency-ban} \\
S17 & Office of the Governor of Oregon. \emph{Governor Kotek issues executive order to advance AI safety and oversight}. 23 Sep 2026. \url{https://apps.oregon.gov/oregon-newsroom/OR/GOV/Posts/Post/governor-kotek-issues-executive-order-to-advance-ai-safety-and-oversight} \\
S18 & OpenAI. \emph{Announcing OpenAI DevDay 2026}. \url{https://openai.com/index/devday-2026/} \\
S19 & TechCrunch. \emph{OpenAI launches GPT-6 Sol and Luna}. 22 Sep 2026. \url{https://techcrunch.com/2026/09/22/openai-launches-gpt-6-sol-and-luna/} \\
\bottomrule
\end{tabularx}}

\vspace{0.8em}
{\footnotesize\color{brandgrey}
Compiled \briefdate. Every figure in this brief was read from the linked primary
source on that date. Corrections: \texttt{\briefauthor}.}

\end{document}
